\clearpage
\fancyhead[L]{\small SONNENDURCHMESSER.}
\hypertarget{native-530-lower}{}\label{native-530-lower}
\begin{center}\small Göttingische gelehrte Anzeigen. Stück 95. Seite 937..940. 1810 Juni 16.\\[3pt]\rule{.72\linewidth}{.3pt}\end{center}

Wir haben im 113. Stücke dieser Gel. Anzeigen vom vorigen Jahre [S. 517 d. B.] die merkwürdigen Resultate mitgetheilt, welche Hr. \textsc{von Lindenau} aus den Greenwicher Beobachtungen von 1750... 1786 über den Sonnendurchmesser gezogen hatte. Seitdem hat dieser Astronom seine Untersuchungen auch auf die neuern Beobachtungen von 1787... 1798 ausgedehnt und die Resultate davon unserm Hrn. Prof. \textsc{Gauss} in einem Schreiben vom 23. Mai mitgetheilt. Es ist interessant, dass dieselben mit den früheren sehr gut harmoniren, und sie verdienen daher um so mehr wenigstens in praktischer Hinsicht alle Aufmerksamkeit. Zuvörderst findet sich auch hier wieder eine Art von periodischer Ungleichheit in den auf die mittlere Entfernung reducirten Halbmessern nach den verschiedenen Jahreszeiten, wie folgende Uebersicht zeigt.

\begin{center} \begin{tabular}{l|c||l|c} & Sonnenhalbmesser & & Sonnenhalbmesser \\ \hline Januar & $959''70$ & Juli & $959''22$ \\ Februar & $959.99$ & August & $959.98$ \\ März & $960.41$ & September & $960.19$ \\ April & $959.80$ & October & $960.10$ \\ Mai & $959.81$ & November & $960.07$ \\ Juni & $959.00$ & December & $958.75$ \\ \end{tabular} \end{center}

Diese Beobachtungen lassen sich ganz gut darstellen, wenn man in der Formel am angeführten Orte S. 1125 [S. 519 d. B.] (wo übrigens durch einen Druckfehler 0.9631 statt 0.09631 steht) $\frac{bb}{aa} - 1 = \frac{1}{92}$ setzt, also etwas kleiner, als der dort bestimmte Werth, jedoch auch positiv.

Die mittlern Sonnenhalbmesser aus allen einzelnen Beobachtungen der verschiedenen Jahre findet Hr. \textsc{von Lindenau}, wie folgt:

\begin{center} \begin{tabular}{l|c|c} & \begin{tabular}{c}Sonnenhalb-\\messer\end{tabular} & \begin{tabular}{c}Anzahl der\\Beobachtungen\end{tabular} \\ \hline 1787 & $959''26$ & 62 \\ 1788 & $959.10$ & 75 \\ 1789 & $959.68$ & 95 \\ 1790 & $959.33$ & 80 \\ 1791 & $959.51$ & 91 \\ 1792 & $960.24$ & 65 \\ 1793 & $959.59$ & 75 \\ 1794 & $959.96$ & 62 \\ 1795 & $960.05$ & 64 \\ 1796 & $960.02$ & 48 \\ 1797 & $959.99$ & 70 \\ 1798 & $960.03$ & 72 \\ \end{tabular} \end{center}

\SourceLeaf{531}
Ueber die scheinbare successive Abnahme des Sonnenhalbmessers gibt die Zusammenstellung der drei und dreissigjährigen \textsc{Maskelyne}'schen Beobachtungen nach drei Epochen folgende Resultate:

\begin{center} \begin{tabular}{c|c} & Sonnenhalbmesser \\ \hline 1765... 1775 & $961''66$ \\ 1776... 1786 & $960.22$ \\ 1787... 1798 & $959.77$ \\ \end{tabular} \end{center}

Aus 116 von \textsc{Piazzi} in den Jahren 1791... 1793 beobachteten horizontalen Sonnenhalbmessern findet Hr. \textsc{von Lindenau} $961''21$, so wie aus \textsc{Bradley}'s Beobachtungen von 1755... 1760, $961''86$. Die nahe Uebereinstimmung dieser Messungen aus so verschiedenen Zeiten bestätigt die Unveränderlichkeit der Grösse des Sonnenkörpers; hingegen von der successiven Abnahme in dem langen Zeitraume der \textsc{Maskelyne}'schen Beobachtungen findet Hr. \textsc{von Lindenau} den Grund in der veränderten Beschaffenheit der Gesichtskraft des Beobachters, wo bei zunehmenden Jahren die Irradiation kleiner geworden zu sein scheint. Hr. \textsc{von Lindenau} erklärt dies selbst nur für eine weiterer Untersuchung nicht unwerthe Vermuthung.

Der ansehnliche Unterschied zwischen den horizontalen und verticalen Halbmessern findet sich auch in den neuern Beobachtungen wieder. Folgendes sind die Resultate der Berechnungen des Hrn. \textsc{von Lindenau} hierüber:

\begin{center}\small\setlength{\tabcolsep}{3pt} \begin{tabular}{l|c|c||l|c|c} & \begin{tabular}{c}Verticalhalb-\\messer\end{tabular} & \begin{tabular}{c}Anzahl der\\Beobachtungen\end{tabular} & & \begin{tabular}{c}Verticalhalb-\\messer\end{tabular} & \begin{tabular}{c}Anzahl der\\Beobachtungen\end{tabular} \\ \hline 1787 & $962''42$ & 72 & 1793 & $963''22$ & 49 \\ 1788 & $962.44$ & 102 & 1794 & $963.13$ & 72 \\ 1789 & $963.10$ & 91 & 1795 & $963.22$ & 60 \\ 1790 & $963.03$ & 66 & 1796 & $962.40$ & 61 \\ 1791 & $963.10$ & 57 & 1797 & $962.25$ & 50 \\ 1792 & $963.30$ & 48 & 1798 & $962.10$ & 53 \\ \end{tabular} \end{center}

Das arithmetische Mittel aus allen ist $962''82$, also um $2''96$ grösser, als das Mittel aus den horizontalen Halbmessern. Aus den Beobachtungen von 1765... 1786 folgte dieser Unterschied $2''80$. Aus 71 \textsc{Piazzi}'schen Beobachtungen findet Hr. \textsc{von Lindenau} den mittlern verticalen Halbmesser $963''01$, und den horizontalen um $1''80$ kleiner. Auch Beobachtungen von \textsc{Bouvard} gaben Hrn. \textsc{von Lindenau} eine ähnliche Differenz.

Sehr auffallend ist diese Uebereinstimmung der Resultate aus den Beobachtungen von drei verschiedenen Beobachtern allerdings. Wenn indess eine länglich-ellipsoidische Gestalt des Sonnenkörpers aus andern überwiegenden Gründen nicht wol für zulässig gehalten werden kann, so möchte man eher geneigt sein, den Grund jenes Unterschiedes anderswo, und vielleicht am füglichsten in der Beobachtungsmanier, zu suchen, in so fern man gewöhnlich bei Messung der Zenithdistanzen den Sonnenrand mit dem Horizontalfaden \emph{von aussen} in Berührung bringt, und daher die Zenithdistanz des obern Sonnenrandes um die halbe Fadendicke zu klein, die des untern um dieselbe Grösse zu gross finden muss; daher denn die auf diese Art bestimmten Sonnendurchmesser um die ganze Fadendicke zu gross ausfallen müssen.
